\section{Introduction}\label{sec:intro}

LLM-driven automated algorithm design (AAD) generates and evolves
optimisation heuristics with a large language model as the variation
operator~\cite{vanstein2025llamea,liu2024eoh,romeraparedes2024funsearch,
ye2024reevo,novikov2025alphaevolve}. The model receives one or more
parent algorithms with their feedback and writes a child, which is
evaluated on a benchmark and selected like any other individual. Methods
differ in their prompts, populations and selection, but all of them
treat the model itself as a black box. This is costly: every generated
algorithm needs a full benchmark run, a sizable fraction of them crash
(16--25\% in our runs, far more on harder problems), and the only way to
change what the model writes is to change the words it reads.

Mechanistic interpretability offers tools to look inside the model. Linear
probes~\cite{alain2016probes} and sparse
autoencoders~\cite{bricken2023monosemanticity,cunningham2024sae}
\emph{read} what is represented in the internal activations, and
activation steering~\cite{turner2023actadd,rimsky2024caa,zou2023repe}
\emph{writes} to them, shifting the behaviour of the model without
changing its prompt. These tools have not been applied to the generator
of an evolutionary loop. For the EC community, they raise four
questions. Does the model know, while it writes, whether its offspring
will work? Which parts of the prompt does it actually use when it
mutates a parent? Does it represent algorithmic concepts in a form we
can find? And can its internal state serve as a cheaper or more
controllable variation operator?

We study these questions with an instrumented open-weight model
(Qwen3-8B, with replications on Qwen3-14B and Gemma-3-12B) inside
\llamea~\cite{vanstein2025llamea}, on MA-BBOB~\cite{vermetten2025mabbob}
and, for transfer, on a photonics problem, circle packing and three BBOB
functions in 20-D, all through the \blade benchmark
suite~\cite{blade2025}. In total we generated and evaluated more than
80,000 algorithms, with their activations, evaluation outcomes and
traced search behaviour. Our contributions are:
\begin{itemize}
  \item \textbf{C1, reading for efficiency (Sec.~\ref{sec:reading}).} The
    child's failure is linearly decodable from the generator's state
    while it writes the code, beyond code text, output uncertainty and
    search metadata. A frozen probe used as a pre-screen skips 21\% of
    the evaluations with non-inferior search quality (pre-registered
    test, 36 seeds), and its ranking transfers to new problems.
  \item \textbf{C2, what drives the edit (Sec.~\ref{sec:reading}).} None of
    three models uses the score written in the feedback; the edits
    follow the parent code and, in the Qwen models, the error trace.
  \item \textbf{C3, concepts (Sec.~\ref{sec:concepts}).}
    Sparse-autoencoder features detect metaheuristic building blocks
    beyond their keywords, and one feature fires on undefined names, a
    common cause of crashes.
  \item \textbf{C4, writing (Sec.~\ref{sec:steering}).} Steering vectors
    as variation operators: concept levers that insert an algorithm
    family, and an exploration dial that shifts the measured search
    behaviour of the generated code and, on five problems it was not
    derived from, the algorithm families written. In the closed loop,
    compared with prompt instructions and a random direction, steering
    towards exploitation is harmful, and a stagnation-adaptive dial is
    the best arm on MA-BBOB (+0.08 \aocc at 200 evaluations; tested on
    new seeds in a pre-registered confirmation).\todo{update with the
    confirmation result}
    A pre-registered landscape test on BBOB 20-D finds no support for
    the intuition that unimodal functions call for exploitation: the
    dial acts on the design space, not on the optimiser's behaviour.
\end{itemize}
Along the way we found sandbox and scoring issues in widely used AAD
benchmarks that inflate failure rates or compress scores
(Sec.~\ref{sec:setup}). Code, prompts and results are available in an
anonymised repository.\todo{anonymised link}
